\documentclass[]{article}

\usepackage{amsmath}
\usepackage{multirow}
\usepackage{natbib}
\usepackage{graphicx} 


\begin{document}

The connection between deterministic chaos in Hamiltonian systems and emergent anomalous transport statistics remains a key challenge, with complex dynamics often simplified to stochastic models like Lévy walks. To probe this relationship, we numerically simulate passive tracer transport in two-dimensional point-vortex systems of varying vortex density ($N$) and contrast the dynamics against a memoryless Lévy walk null model. By analyzing mean squared displacement, velocity autocorrelation, and residence time distributions, while systematically accounting for measurement noise, we uncover a non-monotonic relationship between vortex density and transport behavior. Sparse to intermediate density configurations ($N \le 20$) display near-ballistic motion ($\alpha \approx 2$), whereas the densest system ($N=40$) transitions to a distinct superdiffusive regime ($\alpha \approx 1.68$). Mechanistic analysis reveals this transition is driven by features absent in the stochastic model; the point-vortex system exhibits significant Hamiltonian memory, evidenced by long-lived velocity correlations and heavy-tailed residence time distributions indicative of tracer trapping. Our findings demonstrate that vortex density controls a complex transition between transport regimes and establish that deterministic memory is a crucial feature distinguishing these chaotic flows from their simpler stochastic approximations.

\end{document}
                